\section{Theory and Method}
\label{sec:method}

\subsection{Flow matching and the learned ODE}
\label{sec:flow-matching-background}

In the paired linear-interpolant setting used throughout the low-dimensional and PDE experiments,
a source point $x_0$ and a target point $x_1$ define the conditional path
$x_t=(1-t)x_0+t x_1$ with conditional velocity $u=x_1-x_0$. A standard conditional
flow-matching loss regresses the neural velocity to this conditional velocity at sampled
intermediate points:
\begin{equation}
  \mathcal{L}_{\mathrm{CFM}}(\theta)
  =
  \mathbb{E}_{t,x_0,x_1}
  \left[
    \left\|v_\theta((1-t)x_0+t x_1,t) - (x_1-x_0)\right\|^2
  \right].
  \label{eq:cfm-objective}
\end{equation}
After training, the marginal field $v_\theta$ is integrated by a numerical solver. In this draft
we focus on Euler with a nonuniform grid
\begin{equation}
  0=\tau_0 < \tau_1 < \cdots < \tau_K=1,
  \qquad
  x_{k+1}=x_k+(\tau_{k+1}-\tau_k)v_\theta(x_k,\tau_k).
  \label{eq:euler-grid}
\end{equation}

\subsection{Material derivative and Euler error}
\label{sec:material-error}

Let $x(t)$ be an exact trajectory of \cref{eq:fm-ode}. Consider one step of length
$h_i=\tau_{i+1}-\tau_i$, starting at $x_i=x(\tau_i)$. Write $v_i=v_\theta(x_i,\tau_i)$
for the velocity there and $a_i=\frac{D v_\theta}{D t}(x_i,\tau_i)$ for its rate of change
along the trajectory, as given by \cref{eq:material-derivative}. Expanding the exact solution
about $\tau_i$ gives
\begin{equation}
  x(\tau_{i+1}) = x_i + h_i v_i + \frac{h_i^2}{2}a_i + O(h_i^3).
  \label{eq:euler-expansion}
\end{equation}
Local truncation error is one source of error in the final sample. As integration proceeds,
the learned flow can amplify or damp errors introduced by earlier Euler steps. A further source
is bias in the learned velocity field, which can affect the final sample even when the ODE is
solved accurately. The material derivative tells us where a given Euler step is likely to incur
a larger local error and thus where shorter steps may help.

Starting from the same $x_i$, Euler predicts $x_i+h_i v_i$. Subtracting this prediction
from the exact endpoint in \cref{eq:euler-expansion} gives the local Euler defect
\begin{equation}
  d_i
  :=
  x(\tau_{i+1}) - (x_i+h_i v_i)
  =
  \frac{h_i^2}{2}a_i+O(h_i^3).
  \label{eq:euler-defect}
\end{equation}
For a given step length, $\tfrac12 h_i^2\lVert a_i\rVert$ estimates the magnitude of the
error introduced by that step. To understand how much this error contributes to the final
sample, we must also consider how the subsequent flow amplifies or damps a perturbation
in the direction of the defect.

\subsection{Transporting a local defect to terminal time}
\label{sec:transported-defect}

Let $\Phi_{s,t}$ denote the flow map of the frozen learned ODE from time $t$ to $s$, and define
its state tangent
\begin{equation}
  M(1,t)=D_x\Phi_{1,t}(x(t)).
  \label{eq:flow-tangent}
\end{equation}
Because $d_i$ is injected at the end of Euler step $i$, its leading effect on the terminal state is
\begin{equation}
  d_i^{\mathrm{terminal}}
  \simeq
  M(1,\tau_{i+1})d_i
  \simeq
  \frac{h_i^2}{2}M(1,\tau_{i+1})a_i.
  \label{eq:terminal-defect}
\end{equation}
We therefore distinguish the local and terminally transported profiles
\begin{equation}
  e_i=\lVert a_i\rVert,
  \qquad
  q_i=\left\lVert M(1,\tau_{i+1})a_i\right\rVert.
  \label{eq:transported-profile}
\end{equation}
SCTW uses $e_i$. Adjoint-SCTW uses $q_i$. The latter is closely related to goal-oriented or
dual-weighted residual estimation: a local defect is weighted by the sensitivity of a terminal
quantity of interest \citep{becker2001optimal}.
Neither scalar profile is a complete global-error estimator. In particular, separately taking the
norm of each transported defect discards directional cancellation and nonlinear interactions
between errors injected at different steps. Adjoint-SCTW should therefore be interpreted as a
terminal-sensitivity-weighted mesh heuristic, not as an optimizer of the exact endpoint error.

\subsection{Self-curvature probe}
\label{sec:self-curvature-probe}

SCTW estimates a scalar profile $e(t)$ from the trained model alone. Starting from a small fixed
set of source samples, we run a fine uniform Euler probe with step size $\delta t$. If $x_i$ is
the probe state at time $t_i$ and $x_{i+1}$ is the next state, then
\begin{equation}
  e(t_i)
  \approx
  \frac{
    \left\|
      v_\theta(x_{i+1},t_{i+1}) - v_\theta(x_i,t_i)
    \right\|
  }{\delta t},
  \label{eq:self-curvature-estimator}
\end{equation}
averaged over probe samples. This is a self-curvature estimate: it measures the learned model's
own pathwise velocity variation. It requires no oracle coupling, no pressure field, no
second-order autodifferentiation, and no retraining.

The transported estimators below use the corresponding Heun--Euler defect,
$2(x_{i+1}^{\mathrm{Heun}}-x_{i+1}^{\mathrm{Euler}})/h^2$, along a fine Heun pilot. Both
discretizations are consistent estimators of the same continuum material acceleration, but they
are not numerically identical at finite pilot resolution. We therefore report their normalized
profile discrepancy in the equivalence tier; terminal-sensitivity gains should not be interpreted
as arising solely from this small discretization change.

\subsection{Tangent and finite-difference verification estimators}
\label{sec:transport-estimators}

The action of $M$ on $a_i$ can be evaluated by the tangent-linear model (TLM)
\begin{equation}
  \frac{d\,\delta x}{dt}
  =J_v(x(t),t)\,\delta x,
  \qquad
  \delta x(\tau_{i+1})=a_i,
  \label{eq:tangent-linear-model}
\end{equation}
where $J_v=D_xv_\theta$. Our verification implementation differentiates the same discrete Heun
pilot used for calibration. If one Heun step is
\begin{equation}
  \Psi_j(x)
  =x+\frac{h}{2}\left[
    v_\theta(x,t_j)
    +v_\theta\!\left(x+h v_\theta(x,t_j),t_{j+1}\right)
  \right],
  \label{eq:discrete-heun-step}
\end{equation}
then a Jacobian--vector product (JVP) propagates the perturbation as
\begin{equation}
  \delta x_{j+1}=D\Psi_j(x_j)\,\delta x_j.
  \label{eq:discrete-heun-tangent}
\end{equation}
This is exact for the discrete pilot, up to floating-point and autodifferentiation error. A
centered finite difference provides an independent check along the acceleration direction
$\widehat a_i=a_i/\lVert a_i\rVert$:
\begin{equation}
  q_i
  \approx
  \lVert a_i\rVert
  \frac{\left\lVert
    \Phi_{1,\tau_{i+1}}(x_{i+1}+\varepsilon\widehat a_i)
    -\Phi_{1,\tau_{i+1}}(x_{i+1}-\varepsilon\widehat a_i)
  \right\rVert}{2\varepsilon}.
  \label{eq:transport-finite-difference}
\end{equation}
Finite differences and forward JVPs are useful verification estimators, but each anchor requires a
separate forward propagation. They are therefore not the image-scale default.

\subsection{Randomized discrete adjoint}
\label{sec:randomized-adjoint}

For a terminal probe $z$, the state adjoint at a defect insertion time is
\begin{equation}
  \lambda_{i+1}=M(1,\tau_{i+1})^\top z.
  \label{eq:discrete-adjoint}
\end{equation}
For a zero-mean isotropic probe with $\mathbb{E}[zz^\top]=I$, the transported norm obeys
\begin{equation}
  q_i^2
  =
  \mathbb{E}_z\left[
    \left(a_i^\top\lambda_{i+1}\right)^2
  \right].
  \label{eq:randomized-norm-identity}
\end{equation}
This is the same randomized projection identity that underlies Hutchinson-style matrix estimators
\citep{hutchinson1989stochastic}. With $R$ independent Gaussian terminal probes, we use
\begin{equation}
  \widehat q_i
  =
  \left[
    \frac{1}{R}\sum_{r=1}^{R}
    \left(a_i^\top\lambda_{i+1}^{(r)}\right)^2
  \right]^{1/2}.
  \label{eq:randomized-adjoint-estimator}
\end{equation}
We declare $R=4$ as the scalable default and examine $R\in\{2,4,8\}$ only in the estimator
equivalence tier. The square root makes $\widehat q_i$ a finite-$R$ biased norm estimator even
though its square is unbiased; it is consistent as $R$ grows. Probe-seed and $R$ ablations are
therefore part of the estimator audit rather than hidden tuning. This adjoint is a
state-sensitivity calculation, not parameter-gradient
training: $\theta$ is frozen, no loss is differentiated with respect to model parameters, and no
checkpoint is updated. The connection to adjoint models and neural-ODE adjoints is methodological
rather than a claim that adjoints are new \citep{errico1997adjoint,chen2018neuralode}.

The implementation first runs a no-gradient Heun pilot and stores its states and local
accelerations. During each reverse probe, it recomputes one Heun step, applies its vector--Jacobian
product, and immediately frees that step's graph. Activation memory is therefore independent of
the number of fine pilot steps; no full Jacobian or state-transition matrix is formed.

\subsection{Equidistributed time mesh}
\label{sec:mesh-rule}

For either $r(t)=e(t)$ (SCTW) or $r(t)=q(t)$ (Adjoint-SCTW), if the Euler terminal-error proxy is
$\Delta t_k^2 r(\tau_k)$, then a raw equal-error heuristic would
place steps so that $\Delta t_k \propto r(\tau_k)^{-1/2}$. Equivalently, knots are quantiles of a
density proportional to $\sqrt{r(t)}$. In practice, sharp probe profiles and very small NFE make
this aggressive. We therefore use a tempered density
\begin{equation}
  \rho_p(t) = \left(r(t)+\epsilon\right)^p,
  \qquad p \in \{0.25,0.5\},
  \label{eq:mesh-density}
\end{equation}
and place the $K$ knots by equal cumulative mass:
\begin{equation}
  \int_0^{\tau_j} \rho_p(t)\,dt
  =
  \frac{j}{K}
  \int_0^1 \rho_p(t)\,dt,
  \qquad j=0,\ldots,K.
  \label{eq:equidistribution}
\end{equation}
The exponent $p=0.5$ corresponds to the raw Euler-error rule, while $p=0.25$ is a more conservative
variant that often behaves better at extremely low NFE.

\subsection{Offline population calibration}
\label{sec:offline-calibration}

Both methods use a population-average profile calibrated once from source samples disjoint from
the evaluation priors. The profile, seeds, timing, memory, and all requested NFE grids are cached
and reused for every metric evaluation. Conditional models bind a balanced set of labels to the
velocity during calibration. The calibration cost is reported separately from sampling NFE; when
the same mesh generates $N$ samples, its amortized cost per sample decreases as $1/N$. This design
preserves rectangular GPU batches and differs from a per-sample adaptive ODE solver.

\subsection{Pressure identity as motivation, not implementation}
\label{sec:pressure-identity}

The pressure analysis began as a more ambitious attempt to use the exact structure of linear
interpolant FM. For a linear interpolant, let $u=x_1-x_0$ be the conditional velocity and let
$p(x,t)$ be the marginal density of $x_t$. The marginal velocity is
$v(x,t)=\mathbb{E}[u \mid x_t=x]$ and the conditional covariance is
\begin{equation}
  \Sigma(x,t)
  =
  \operatorname{Cov}[u \mid x_t=x].
  \label{eq:conditional-covariance}
\end{equation}
For this marginal field, the material derivative satisfies the exact momentum-balance identity
\begin{equation}
  \frac{D v}{D t}
  =
  -\frac{1}{p}\nabla_x \cdot \left(p\,\Sigma\right)
  =
  -\nabla_x \cdot \Sigma - \Sigma \nabla_x \log p.
  \label{eq:pressure-identity}
\end{equation}
This identity explains why the marginal velocity can be curved even when every conditional path
is a straight line: superposed conditional trajectories produce a pressure tensor. It also explains
why forcing zero acceleration can erase meaningful transport structure. The deployed SCTW sampler
does not compute \cref{eq:pressure-identity}; it uses the self-probe in
\cref{eq:self-curvature-estimator}. The pressure identity is the theoretical lens that led us to
measure material derivative as the local integration difficulty.
